\section{ضریب همدوسی}
ضریب همدوسی متقابل برای دو ماتریس متعامد و یکّه‌ی
$U_{m\times m}$
و
$V_{m\times m}$
$\big(U^HU=V^HV=I_{m\times m}\big)$
به صورت زیر تعریف می‌شود:
\begin{equation}
\nonumber
\mu(U,V)\eqdef \max_{1\leq i,j\leq m}\big\lvert\langle u_i,v_j\rangle\big\rvert
\end{equation}

که در آن 
$u_i$
و
$v_j$
، به ترتیب نمایشگر ستون
$i$
و
$j$
از ماتریس‌های
$U$
و
$V$
است.

\begin{itemize}
\item
  ثابت کنید:
\begin{align}
\nonumber\frac{{1}}{\sqrt{m}}\leq \mu(U,V)\leq {1}
\end{align}
\item
 اکنون فرض‌کنید درایه‌های این دو ماتریس دارای توزیع 
$i.i.d$
برنولی به صورت زیر  باشند:

\[
u_{ij},v_{ij} = 
\begin{cases} 
{1/\sqrt{m}} &\quad \frac{1}{2} \text{با احتمال }\\
-{1/\sqrt{m}} &\quad \frac{1}{2} \text{با احتمال }
\end{cases}
\] 

در این صورت نشان دهید که با احتمال حداقل 
$1-\epsilon^{2}$
برای ضریب همدوسی متقابل خواهیم داشت:
\begin{equation}
\nonumber
\mu(U,V)\leq {2\sqrt{\frac{\ln{(\sqrt{2}m/\epsilon)}}{m}}}
\end{equation}
\end{itemize}